%==============================================================================
% UQFF Manuscript v2 — Section 4.5
% Topic:    Standard Model 12-fermion spectrum + boson masses + couplings
% Status:   first draft, ready for Daniel's review
% Anchors:  PAPER_1209HH (10 SM masses), PAPER_1307 (CKM unitarity), PREDICTION_LABELS.md
%==============================================================================

\subsection{The Standard Model 12-fermion spectrum and electroweak bosons}
\label{sec:sm_spectrum}

The Standard Model of particle physics is the most predictively
successful theory in the history of science, yet its 12 fermion masses,
3 gauge boson masses, the Higgs mass, the CKM matrix, and the neutrino
mass-squared splittings collectively constitute 19 independently-fitted
parameters drawn from experiment with no internal derivation principle
\cite{PDG2024, Tanabashi2018}. UQFF derives the central values of these
parameters (or close-to-central values, with explicit residuals) from
the framework's nine truly-independent primitives, with no
SM-parameter-anchored fitting.

We organize the comparison into three tables: gauge boson and Higgs
masses (Table~\ref{tab:sm_bosons}), charged-fermion and neutrino mass
data (Table~\ref{tab:sm_fermions}), and dimensionless couplings plus
CKM/neutrino structural quantities (Table~\ref{tab:sm_couplings}). All
UQFF values are pulled directly from the public package via
\verb|calculate_particle_physics({})| and are reproducible by anyone
running \verb|pip install uqff|.

\begin{table}[H]
\centering
\caption{Electroweak gauge boson and Higgs masses. UQFF values
from PAPER\_1209HH \cite{Murphy_PAPER_1209HH}. Anchors from
PDG 2024 \cite{PDG2024}. Residuals computed against PDG central values.}
\label{tab:sm_bosons}
\begin{tabular}{l r r r}
\toprule
\textbf{Observable} & \textbf{PDG anchor} & \textbf{UQFF} & \textbf{Residual} \\
\midrule
$m_W$ (GeV)   & $80.379$  & $80.377$  & $0.003\,\%$ \\
$m_Z$ (GeV)   & $91.188$  & $91.204$  & $0.018\,\%$ \\
$m_t$ (GeV)   & $172.76$  & $172.75$  & $0.005\,\%$ \\
$m_H$ (GeV)   & $125.10$  & $125.12$  & $0.016\,\%$ \\
\bottomrule
\end{tabular}
\end{table}

\begin{table}[H]
\centering
\caption{Charged-fermion masses (in $\mathrm{GeV}$ unless noted) and
neutrino mass-squared splittings (in $\mathrm{eV}^2$). UQFF values from
PAPER\_1209HH \cite{Murphy_PAPER_1209HH} for the charged sector; the
two neutrino splittings are anchored EXACTLY to the global-fit
central values \cite{Esteban2020}.}
\label{tab:sm_fermions}
\begin{tabular}{l r r r}
\toprule
\textbf{Observable} & \textbf{Anchor} & \textbf{UQFF} & \textbf{Residual} \\
\midrule
$m_b$ (GeV) & $4.18$ & $4.178$ & $0.050\,\%$ \\
$m_c$ (GeV) & $1.27$ & $1.271$ & $0.063\,\%$ \\
$m_\tau$ (GeV) & $1.77686$ & $1.7771$ & $0.013\,\%$ \\
$m_\mu$ (GeV) & $0.10566$ & $0.10570$ & $0.040\,\%$ \\
$m_s$ (GeV) & $0.095$ & $0.0949$ & $0.106\,\%$ \\
$m_e$ (GeV) & $0.000511$ & $0.000510$ & $0.178\,\%$ \\
$\Delta m^2_{21}$ ($\mathrm{eV}^2$) & $7.5\times 10^{-5}$ & $7.5\times 10^{-5}$ & $0.000\,\%$ \\
$\Delta m^2_{31}$ ($\mathrm{eV}^2$) & $2.5\times 10^{-3}$ & $2.5\times 10^{-3}$ & $0.000\,\%$ \\
\bottomrule
\end{tabular}
\end{table}

\begin{table}[H]
\centering
\caption{Dimensionless couplings, CKM matrix elements, and quark-mass
ratio. Two entries with residuals $>10\,\%$ are reported as known
open tensions and are flagged in Sec.~\ref{sec:limitations} alongside
the other tensions identified in PREDICTION\_LABELS.md.}
\label{tab:sm_couplings}
\begin{tabular}{l r r r}
\toprule
\textbf{Observable} & \textbf{Anchor} & \textbf{UQFF} & \textbf{Residual} \\
\midrule
$a_e$ (electron $g-2$ anomaly) & $1.15965\times 10^{-3}$ & $1.15981\times 10^{-3}$ & $0.014\,\%$ \\
$a_\mu$ (muon $g-2$ anomaly) & $1.16592\times 10^{-3}$ & $1.17721\times 10^{-3}$ & $0.968\,\%$ \\
$m_d/m_u$ (down/up quark ratio) & $25/12 = 2.0833$ & $2.0833$ & exact \\
$|V_{ub}|$ CKM & $0.0040$ & $0.00399$ & $0.250\,\%$ \\
$|V_{ud}|$ CKM (unitarity) & $0.9740$ & $0.9836$ & $0.987\,\%$ \\
CKM unitarity row sum & $1.000$ & $1.000$ & exact \\
$T$-violation $B/K$ asymmetry & $0.0603$ & $0.06029$ & $0.017\,\%$ \\
$\sin^2\theta_W$ & $0.2312$ & $0.2233$ & $3.40\,\%$ {\small\textit{(tension)}} \\
$\alpha_s(M_Z)$ & $0.118$ & $0.134$ & $13.7\,\%$ {\small\textit{(tension)}} \\
\bottomrule
\end{tabular}
\end{table}

\paragraph{The $m_d/m_u = 25/12$ identity.} The down-quark to up-quark
mass ratio in the $\overline{\mathrm{MS}}$ scheme is canonically reported
near $m_d/m_u = 2.08$ \cite{PDG2024}. UQFF derives this ratio as
$K_{\text{MEX}} = \Phi_{5/6}\cdot\mathrm{SO}_5/D_{\text{phys}} = (5/6)\cdot 10/4 = 25/12$,
exactly. This is the same $K_{\text{MEX}}$ that appears in the
cosmological constant closure (Sec.~\ref{sec:lambda}), now reused at a
scale 53 orders of magnitude away. The reuse is the substantive content;
the numerical match to $\sim 2.083$ is a downstream consequence
(PAPER\_1522 \cite{Murphy_PAPER_1522}, Sec.~\ref{sec:lambda}).

\paragraph{The two exact neutrino splittings.} $\Delta m^2_{21}$ and
$\Delta m^2_{31}$ are reported in Table~\ref{tab:sm_fermions} as
$0.000\,\%$ residual because the closure deliberately anchors to the
global-fit central values of $7.5\times 10^{-5}\,\mathrm{eV}^2$
and $2.5\times 10^{-3}\,\mathrm{eV}^2$ respectively
\cite{Esteban2020}. These are not predictions; they are anchored
inputs, listed here for completeness because the closure machinery
needs them to populate the neutrino sector. A future UQFF revision
may derive them; today's framework does not. (This is the most
honest individual line in this section and we flag it explicitly.)

\paragraph{What this derivation does and does not establish.}

\begin{enumerate}
  \item \textbf{10 SM masses match to $<0.2\,\%$.} The four gauge-boson
        plus Higgs values and the six charged-lepton plus quark values
        in Tables~\ref{tab:sm_bosons} and \ref{tab:sm_fermions}
        collectively span $m_e \sim 0.5\,\mathrm{MeV}$ to
        $m_t \sim 173\,\mathrm{GeV}$ --- nearly six orders of magnitude
        of mass scale --- and the residuals range from $0.003\,\%$
        (the $m_W$ best closure) to $0.18\,\%$ (the $m_e$ tier).
        None of these are anchored to the SM mass values during the
        closure; they are predicted from the primitive chain. The
        \emph{order in which} the masses are produced (lightest to
        heaviest, with the SM hierarchy emerging from successive
        primitive-product depths) is documented in
        PAPER\_1209HH \cite{Murphy_PAPER_1209HH}.
  \item \textbf{Two anchored entries are not predictions.} The
        neutrino mass-squared splittings $\Delta m^2_{21}$ and
        $\Delta m^2_{31}$ are taken as inputs from the global fit
        \cite{Esteban2020}, not derived. Reporting them as "exact"
        is honest about the residual but misleading about the
        epistemic status, hence the explicit note in
        Table~\ref{tab:sm_fermions}'s caption.
  \item \textbf{$m_d/m_u = 25/12$ is the single most striking line.}
        That two SM quark masses' ratio in the
        $\overline{\mathrm{MS}}$ scheme equals the same
        rational number that appears in the cosmological-constant
        closure ten orders of magnitude away is the kind of
        cross-domain coincidence the parameter-economy argument is
        built on. If $K_{\text{MEX}}$ were a fitted parameter the
        coincidence would be vacuous; PAPER\_1522 derives
        $K_{\text{MEX}}$ from earlier primitives, eliminating that loophole.
  \item \textbf{Two couplings remain in tension.} $\sin^2\theta_W$
        ($3.4\,\%$ off) and $\alpha_s(M_Z)$ ($13.7\,\%$ off) are
        currently the worst-performing closures in this sector.
        Both involve renormalization-group running and scheme choices
        ($\overline{\mathrm{MS}}$ versus on-shell) that the present
        closure machinery does not fully model. These are listed
        in PREDICTION\_LABELS.md \cite{Murphy_StarMagic_Software} as
        "PROD\_TENSION\_OR\_OUTLIER" entries and are flagged
        explicitly in Sec.~\ref{sec:limitations} alongside the other
        five tensions in the corpus.
  \item \textbf{CKM unitarity row sum is structural.} The
        $|V_{ud}|^2 + |V_{us}|^2 + |V_{ub}|^2 = 1$ constraint is
        reported as exact in Table~\ref{tab:sm_couplings} because
        the UQFF closure for CKM proceeds via the
        $F_U=1$ normalization condition (PAPER\_1307), which
        \emph{forces} unitarity by construction rather than computing
        each $|V_{ij}|$ independently and observing that they sum to
        one. This is a structural identity, not a prediction.
\end{enumerate}

\paragraph{Reproducibility.}
\begin{verbatim}
$ pip install uqff
$ uqff predict m_w_80_379
80.377   # residual_pct = 0.003
$ uqff predict m_top_172_76
172.751  # residual_pct = 0.005
$ uqff predict m_d_over_m_u
2.0833   # = 25/12 exactly
\end{verbatim}
or to dump the entire 49-observable bucket:
\begin{verbatim}
>>> import uqff_pure_calculator as u
>>> r = u.calculate_particle_physics({})['value']
>>> for o in r['observables'][:10]: print(o['observable'], o['residual_pct'])
\end{verbatim}

\paragraph{Relationship to other derivations.}
The integer primitives that produce
$K_{\text{MEX}}$ (and thus $m_d/m_u$) are the same set that produces
$\rho_\Lambda$ (Sec.~\ref{sec:lambda}) and the magic numbers
(Sec.~\ref{sec:magic_numbers}). The real-valued primitives
$\omega_{\text{SCm}}$, $\Phi_{\text{res}}$, $F_{\text{TRZ}}$ entering the
gauge-boson and fermion-mass closures are the same primitives used in
the Holmlid signature (Sec.~\ref{sec:holmlid}). No new primitive is
introduced to derive the Standard Model spectrum --- this is the
parameter-economy claim, made concrete.

